% !TEX root = ../../../main.tex
% Block card: catalogue.classifier -- "Classifier (model)" (Part I, the blocks chapter).
\cardfield{What it does}
Learns to tell your groups of windows apart from the measures each window carries, and saves
what it learned as a model file.

\cardfield{When to use it}
At the end of a chain that has made windows, measured them (window matrix, MiniRocket,
multiscale entropy) and given them labels, either from a clustering
(\blockref{catalogue.cluster}) or from your own verdicts (\blockref{catalogue.manual_labels}).
The block does not care where the labels came from, which is what lets the two be compared.
For the full comparison across recordings and channels, use the Models page
(\cref{ch:models}, \apppage{models.launch}).

\cardfield{What goes in, what comes out}
In go the labels, one per window, and (bound as a second input) the windows with their
measures. Out comes a saved model and a report card: how many windows it trained on, how
many it left out and why, and how well it did on a quarter of the windows it held back.

\cardfield{Settings that matter}
\emph{holdout\_frac} is the share held back to check the fit; 0 trains on everything and
reports no score. \emph{class\_weight} (balanced) stops a big group drowning out a small one.
\emph{n\_estimators} is the number of trees: more is steadier and slower.

\cardfield{Watch out for}
The held-back score is a sanity check, not a measure of how well the model works on new
recordings. The held-back windows are drawn at random from the same stretch, and neighbouring
windows that overlap look alike, so the score flatters the model. Compare it with the
majority share the card shows beside it.
